% ----------------------------------------------------------------------
%  3.1.2. Электронное строение атомов и периодическая система
% ----------------------------------------------------------------------
\subsection{Электронное строение атомов и периодическая система}\label{ssec:atoms}

\begin{orientir}
Свойства твёрдого тела определяются внешними (валентными) электронами
атомов. Чтобы понять, сколько их и как они расположены, нужна квантовая
механика. Здесь кратко повторяются её основные результаты для атома
водорода, затем~— для многоэлектронных атомов, и в конце~— как из
строения оболочек вытекает периодичность свойств элементов (размеров
атомов, энергий ионизации, температур плавления).
\end{orientir}

\subsubsection{Волны и частицы}

Свет в явлениях интерференции и дифракции ведёт себя как волна с
частотой $\nu$, а в фотоэффекте и эффекте Комптона~— как поток частиц
(фотонов) с энергией $E=h\nu$ и импульсом $p=h/\lambda$. Де Бройль
предположил, что такая двойственность универсальна: любая частица
массы $m$, движущаяся со скоростью $\upsilon$, ведёт себя и как волна с
длиной
\begin{keyeq}\label{eq:debroglie}
\lambda=\frac hp=\frac h{m\upsilon} .
\end{keyeq}

\begin{primer}[Длина волны электрона]
Электрон, ускоренный напряжением 100~В, имеет энергию $E=100\eV$ и
импульс $p=\sqrt{2m_eE}$, поэтому
$\lambda=h/\sqrt{2m_eE}\approx1{,}23/\sqrt{100}~\text{нм}=0{,}123$~нм~—
порядка межатомных расстояний. Поэтому электроны (а также нейтроны)
дифрагируют на кристаллах; на этом основаны электронография и
нейтронография.
\end{primer}

\subsubsection{Уравнение Шрёдингера и волновая функция}

Движение микрочастиц описывает квантовая механика. Если потенциальная
энергия $U$ не зависит от времени, состояние частицы с энергией $E$
описывается \term{стационарным уравнением Шрёдингера}
\begin{keyeq}\label{eq:schrodinger}
-\frac{\hbar^2}{2m}\Delta\psi+U\psi=E\psi
\qquad\Longleftrightarrow\qquad
\Delta\psi+\frac{8\pi^2m}{h^2}(E-U)\psi=0 ,
\end{keyeq}
где $\hbar=h/2\pi$, $\Delta$~— оператор Лапласа. Уравнение не выводится
из более общих законов: оно постулируется (как уравнение стоячей волны с
длиной волны де Бройля) и проверяется опытом. Точно оно решается только
для одноэлектронных систем (H, He$^+$, Li$^{2+}$, \ldots), где не нужно
учитывать взаимодействие электронов.

\begin{opredelenie}
\term{Волновая функция} $\psi$ сама по себе непосредственного физического
смысла не имеет. Смысл имеет квадрат её модуля: $|\psi|^2\,\d V=\psi\psi^*\d V$~—
\emph{вероятность} найти частицу в элементе объёма $\d V$; $|\psi|^2$~—
\term{плотность вероятности}.
\end{opredelenie}

Чтобы решение имело смысл, $\psi$ должна быть конечной, непрерывной,
однозначной, иметь непрерывную первую производную, обращаться в нуль
там, где частицы быть не может (например, на бесконечности), и
удовлетворять \term{условию нормировки} $\int\psi\psi^*\d V=1$
(где-нибудь частица находится наверняка).

Таким условиям решения удовлетворяют лишь при определённых значениях
энергии~— \term{собственных значениях}; соответствующие решения~—
\term{собственные функции}. Так квантование энергии получается
автоматически.

\subsubsection{Атом водорода: квантовые числа}

Решение уравнения Шрёдингера для атома водорода даёт три главных
результата.

\begin{enumerate}[label=\arabic*)]
\item \emph{Энергия квантуется} (так же, как в теории Бора):
\begin{keyeq}\label{eq:H-levels}
E_n=-\frac{2\pi^2me^4}{h^2}\cdot\frac1{n^2}=-\frac{13{,}6\eV}{n^2},\qquad n=1,2,3,\ldots
\end{keyeq}
$n$~— \term{главное квантовое число}.
\item \emph{Орбитальный момент импульса квантуется}:
$M_l=\hbar\sqrt{l(l+1)}$, где $l=0,1,\ldots,n-1$~— \term{орбитальное
квантовое число}.
\item \emph{Проекция момента} на выделенное направление (например, на
направление магнитного поля) равна $m_l\hbar$, где
$m_l=-l,\ldots,0,\ldots,l$ ($2l+1$ значений)~— \term{магнитное квантовое
число}.
\end{enumerate}
Каждому набору $(n,l,m_l)$ отвечает своё состояние $\psi_{nlm_l}$. При
данном $n$ их $\sum_{l=0}^{n-1}(2l+1)=n^2$. Состояния с $l=0,1,2,3$
обозначают буквами $s$, $p$, $d$, $f$ и пишут перед буквой значение $n$:
$1s$, $2p$, $3d$\ldots\ (рис.~\ref{fig:H-cells}).

\begin{figure}[ht]
\centering
\begin{tikzpicture}[scale=0.85]
  \def\bw{0.42}
  \draw[arr] (-0.6,-0.4) -- (-0.6,5.4) node[above, slbl]{$E$};
  \foreach \n/\y/\E in {1/0/{-13{,}6}, 2/2.8/{-3{,}4}, 3/4.0/{-1{,}5}, 4/4.7/{-0{,}85}}{
    \draw[fmGray, dashed] (-0.6,\y) -- (11.4,\y);
    \node[slbl, left] at (-0.7,\y) {$n=\n$};
    \node[slbl, right] at (11.4,\y) {$\E$~эВ};
  }
  % ячейки: s (1), p (3), d (5), f (7)
  \newcommand{\cells}[3]{\foreach \k in {1,...,#3}{\draw[fill=fmBlue!10] ({#1+(\k-1)*\bw},{#2-0.2}) rectangle ++(\bw,0.4);}}
  \cells{0}{0}{1}
  \cells{0}{2.8}{1}\cells{1.2}{2.8}{3}
  \cells{0}{4.0}{1}\cells{1.2}{4.0}{3}\cells{3.2}{4.0}{5}
  \cells{0}{4.7}{1}\cells{1.2}{4.7}{3}\cells{3.2}{4.7}{5}\cells{6.2}{4.7}{7}
  \node[slbl] at (0.2,5.3) {$s$}; \node[slbl] at (1.83,5.3) {$p$};
  \node[slbl] at (4.25,5.3) {$d$}; \node[slbl] at (7.67,5.3) {$f$};
  \node[slbl] at (0.2,-0.55) {$l=0$}; \node[slbl] at (1.83,2.25) {$l=1$};
  \node[slbl] at (4.25,3.45) {$l=2$};
\end{tikzpicture}
\caption{Квантовые ячейки атома водорода: каждая ячейка~— состояние с
определёнными $n$, $l$, $m_l$; в ячейке могут находиться два электрона с
противоположными спинами. В атоме водорода энергия зависит только от $n$
(масштаб по вертикали условный)}
\label{fig:H-cells}
\end{figure}

\subsubsection{Форма электронного облака}

Траекторию электрона указать нельзя (принцип неопределённости
Гейзенберга), но можно найти распределение вероятности~— «электронное
облако». В сферических координатах $(r,\vartheta,\varphi)$ волновая
функция атома водорода распадается на произведение множителей:
радиального $R_{nl}(r)$ и углового $A_{lm}(\vartheta,\varphi)$. Поэтому
форму облака удобно описывать по отдельности: сколько электронной
плотности находится на расстоянии $r$ от ядра и как она распределена по
направлениям.

\paragraph{Радиальное распределение}
Удобно рассматривать не саму плотность $|\psi|^2$, а
\term{радиальную плотность} $U(r)=4\pi r^2|\psi|^2$~— вероятность найти
электрон в тонком шаровом слое радиусом $r$ и единичной толщины
(рис.~\ref{fig:radial}). Например,
\[
\psi_{1s}=\frac1{\sqrt{\pi a^3}}\,e^{-r/a},\qquad
\psi_{2s}=\frac1{2\sqrt{2\pi a^3}}\left(1-\frac r{2a}\right)e^{-r/2a},
\]
где $a=0{,}0529$~нм~— радиус первой боровской орбиты. Из условия
$\d U/\d r=0$ максимум $U(r)$ для $1s$-состояния приходится ровно на
$r=a$~— в согласии с теорией Бора.

\begin{figure}[ht]
\centering
\begin{subfigure}[b]{0.49\textwidth}\centering
\begin{tikzpicture}
\begin{axis}[width=\linewidth, height=5cm, xmin=0, xmax=22, ymin=0, ymax=0.56,
  xlabel={$r/a$}, ylabel={$U(r)\cdot a$}, legend style={font=\scriptsize}, legend cell align=left]
  \addplot[fmBlue, very thick, domain=0:22, samples=200] {4*x^2*exp(-2*x)};
  \addlegendentry{$1s$}
  \addplot[fmGreen, very thick, domain=0:22, samples=200] {0.5*x^2*(1-x/2)^2*exp(-x)};
  \addlegendentry{$2s$}
  \addplot[fmRed, very thick, domain=0:22, samples=250] {(4/27)*x^2*(1-2*x/3+2*x^2/27)^2*exp(-2*x/3)};
  \addlegendentry{$3s$}
\end{axis}
\end{tikzpicture}
\caption{}
\end{subfigure}\hfill
\begin{subfigure}[b]{0.49\textwidth}\centering
\begin{tikzpicture}
\begin{axis}[width=\linewidth, height=5cm, xmin=0, xmax=22, ymin=0, ymax=0.25,
  xlabel={$r/a$}, ylabel={$U(r)\cdot a$}, legend style={font=\scriptsize}, legend cell align=left]
  \addplot[fmGreen, very thick, domain=0:22, samples=200] {(1/24)*x^4*exp(-x)};
  \addlegendentry{$2p$}
  \addplot[fmRed, very thick, domain=0:22, samples=200] {(64/4374)*x^4*(1-x/6)^2*exp(-2*x/3)};
  \addlegendentry{$3p$}
\end{axis}
\end{tikzpicture}
\caption{}
\end{subfigure}
\caption{Радиальная плотность $U(r)=4\pi r^2|\psi|^2$ для $s$-состояний
(\textit{а}) и $p$-состояний (\textit{б}) атома водорода. Внешние
максимумы приходятся примерно на радиусы боровских орбит ($a$, $4a$,
$9a$); у $2s$ и $3s$ есть сферические узлы~— расстояния, где плотность
обращается в нуль}
\label{fig:radial}
\end{figure}

Видно, что у $1s$-состояния плотность обращается в нуль только на
бесконечности, а у $2s$ и $3s$~— ещё и на определённых расстояниях
(\term{сферические узлы}); у $ns$-состояния их $n-1$, у $np$~— $n-2$
(и узел в центре). Строго говоря, при учёте спина (релятивистская теория
Дирака) плотность в узлах не равна нулю, а очень мала.

\paragraph{Угловое распределение}
У $s$-состояний ($l=0$) угловая часть постоянна: облако сферически
симметрично. У $p$-состояний ($l=1$) угловые множители удобно выбрать
действительными, пропорциональными $\cos\vartheta$,
$\sin\vartheta\cos\varphi$ и $\sin\vartheta\sin\varphi$, т.\,е.
вытянутыми вдоль осей $z$, $x$ и $y$ (рис.~\ref{fig:orbitals}). Эти три
«гантели» вместе снова дают сферически симметричное облако. У $d$- и
$f$-состояний форма сложнее (например, четырёхлепестковая у $d_{xy}$).

\begin{figure}[ht]
\centering
\begin{tikzpicture}[scale=1.0]
  \tikzset{lobe/.style={fill=#1!35, draw=#1!80!black, line width=0.5pt}}
  \begin{scope}
    \fill[fmBlue!30, draw=fmBlue] (0,0) circle (0.8);
    \node[slbl] at (0,-1.4) {$s$};
  \end{scope}
  \begin{scope}[shift={(3,0)}]
    \draw[fmGray!70] (-1.2,0) -- (1.2,0) (0,-1.2) -- (0,1.2);
    \draw[lobe=fmGreen] plot[domain=-90:90, samples=60, smooth] ({1.1*cos(\x)^2*cos(\x)},{1.1*cos(\x)^2*sin(\x)}) -- cycle;
    \draw[lobe=fmGreen] plot[domain=90:270, samples=60, smooth] ({1.1*cos(\x)^2*cos(\x)},{1.1*cos(\x)^2*sin(\x)}) -- cycle;
    \node[slbl] at (1.35,0.15) {$x$};
    \node[slbl] at (0,-1.4) {$p_x$};
  \end{scope}
  \begin{scope}[shift={(6,0)}]
    \draw[fmGray!70] (-1.2,0) -- (1.2,0) (0,-1.2) -- (0,1.2);
    \draw[lobe=fmRed] plot[domain=0:180, samples=60, smooth] ({1.1*sin(\x)^2*cos(\x)},{1.1*sin(\x)^2*sin(\x)}) -- cycle;
    \draw[lobe=fmRed] plot[domain=180:360, samples=60, smooth] ({1.1*sin(\x)^2*cos(\x)},{1.1*sin(\x)^2*sin(\x)}) -- cycle;
    \node[slbl] at (0.2,1.35) {$z$};
    \node[slbl] at (0,-1.4) {$p_z$};
  \end{scope}
  \begin{scope}[shift={(9,0)}]
    \draw[fmGray!70] (-1.2,0) -- (1.2,0) (0,-1.2) -- (0,1.2);
    \foreach \q in {0,90,180,270}{
      \draw[lobe=fmAmber] plot[domain=\q:\q+90, samples=60, smooth] ({1.1*sin(2*\x)^2*cos(\x)},{1.1*sin(2*\x)^2*sin(\x)}) -- cycle;}
    \node[slbl] at (0,-1.4) {$d_{xy}$};
  \end{scope}
\end{tikzpicture}
\caption{Угловое распределение электронной плотности (сечение плоскостью
рисунка): $s$~— сфера, $p_x$ и $p_z$~— «гантели» вдоль осей, $d_{xy}$~—
четыре лепестка между осями}
\label{fig:orbitals}
\end{figure}

\paragraph{Спин}
Кроме орбитального движения, у электрона есть собственный момент
импульса~— \term{спин}. Его проекция на выделенное направление
принимает лишь два значения $m_s\hbar$, $m_s=\pm\frac12$~—
\term{спиновое квантовое число}. Итак, состояние электрона в атоме
задаётся четырьмя квантовыми числами $(n,l,m_l,m_s)$.

\subsubsection{Многоэлектронные атомы}

Для атомов с несколькими электронами уравнение Шрёдингера точно не
решается: мешает взаимодействие электронов друг с другом. Приближённые
методы (вариационный, самосогласованного поля и др.) дают главный
вывод: \emph{электронная оболочка многоэлектронного атома состоит из
слоёв, подобных уровням атома водорода}.

\begin{itemize}
\item Слои: $n=1$~— $K$-слой, $n=2$~— $L$, $n=3$~— $M$, $n=4$~— $N$\ldots
\item \term{Принцип Паули}: в одной квантовой ячейке $(n,l,m_l)$ могут
находиться не более двух электронов, причём с противоположными спинами.
Отсюда ёмкость оболочки $(n,l)$~— $2(2l+1)$ электронов ($s$~— 2, $p$~— 6,
$d$~— 10, $f$~— 14), ёмкость слоя~— $2n^2$.
\item В отличие от водорода, энергия зависит не только от $n$, но и от $l$:
она растёт с ростом обоих чисел. Чем больше электронов в атоме, тем
сильнее зависимость от $l$, и при $n\ge4$ уровни $(n-1)d$ и $(n-2)f$
оказываются близки к уровню $ns$ и даже ниже. Это определяет особенности
переходных металлов, лантаноидов и актиноидов.
\end{itemize}

\paragraph{Опытное доказательство слоистого строения: потенциалы ионизации}
\term{Потенциал ионизации} $J_k$ $k$-го порядка~— работа отрыва $k$-го
электрона от атома, у которого уже удалены $k-1$ электронов. Таблица
\ref{tab:ionization} показывает потенциалы ионизации первых элементов
(в единицах потенциала ионизации водорода $13{,}6\eV$).

\begin{table}[ht]
\centering
\caption{Потенциалы ионизации первых элементов (в единицах $J_{\text H}=13{,}6\eV$).
Жирным выделены скачки~— начало отрыва электронов внутреннего слоя $1s$}
\label{tab:ionization}
\small
\begin{tabular}{@{}llccccccc@{}}
\toprule
Элемент & Оболочки & $J_1$ & $J_2$ & $J_3$ & $J_4$ & $J_5$ & $J_6$ & $J_7$ \\
\midrule
H  & $1s$         & 1,00 \\
He & $1s^2$       & 1,81 & 4,00 \\
Li & $1s^22s$     & 0,40 & \textbf{5,56} & 9,01 \\
Be & $1s^22s^2$   & 0,69 & 1,34 & \textbf{11,3} & 16,0 \\
B  & $1s^22s^22p$ & 0,61 & 1,85 & 2,79 & \textbf{19,1} & 25,0 \\
C  & $1s^22s^22p^2$ & 0,83 & 1,79 & 3,52 & 4,74 & \textbf{28,8} & 36,0 \\
N  & $1s^22s^22p^3$ & 1,07 & 2,18 & 3,49 & 5,70 & 7,20 & \textbf{40,6} & 49,1 \\
\bottomrule
\end{tabular}
\end{table}

Что видно из таблицы\ispr{В табл.~3.1 оригинала часть значений искажена
(например, для углерода $J_3=8{,}52$ вместо $3{,}52$, для бора $J_2=1{,}65$
вместо $1{,}85$, для азота $J_4=5{,}43$ вместо $5{,}70$). Здесь приведены
справочные значения.}:
\begin{itemize}
\item $J_1$ гелия почти вдвое больше, чем у водорода: заряд ядра вдвое
больше, а второй электрон лишь немного экранирует ядро.
\item У лития $J_1$ резко падает (в 4,5 раза): отрывается электрон
$2s$~— он во внешнем слое и слабо связан. Зато $J_2$ лития в 14 раз
больше $J_1$: второй электрон уже из слоя $n=1$. Такие скачки (выделены
жирным) отмечают границы слоёв.
\item Внутри слоя последовательные потенциалы растут примерно
пропорционально заряду иона.
\item Последние потенциалы ионизации (одноэлектронные ионы) относятся
как квадраты зарядов ядер, как и предсказывает теория
водородоподобных атомов: $J\propto Z^2$.
\end{itemize}

\paragraph{Рентгеновские спектры}
Ещё одно доказательство~— рентгеновские спектры. При уменьшении длины
волны (росте энергии кванта) коэффициент поглощения $\mu$ плавно
убывает ($\mu\propto\lambda^3$), но как только энергия кванта становится
достаточной, чтобы выбить электрон из $L$- или $K$-слоя, $\mu$ скачком
возрастает (рис.~\ref{fig:xray},\,а)~— это \term{края поглощения}. А при
переходах электронов из внешних слоёв на освободившиеся места во
внутренних испускается \term{характеристическое излучение}: серии $K$,
$L$, $M$ (рис.~\ref{fig:xray},\,б).

\begin{figure}[ht]
\centering
\begin{subfigure}[b]{0.48\textwidth}\centering
\begin{tikzpicture}
\begin{axis}[width=\linewidth, height=4.8cm, ymode=log, xmin=0.2, xmax=4.5, ymin=0.05, ymax=200,
  xlabel={$\lambda$}, ylabel={$\mu$}, xtick={1,3}, xticklabels={$\lambda_K$,$\lambda_L$}, ytick=\empty]
  \addplot[fmBlue, very thick, domain=0.2:1, samples=40] {8*x^3};
  \addplot[fmBlue, very thick, domain=1:3, samples=40] {1.0*x^3};
  \addplot[fmBlue, very thick, domain=3:4.5, samples=40] {0.15*x^3};
  \draw[fmBlue, very thick] (axis cs:1,8) -- (axis cs:1,1);
  \draw[fmBlue, very thick] (axis cs:3,27) -- (axis cs:3,4.05);
  \node[font=\footnotesize] at (axis cs:1.3,20) {$K$-край};
  \node[font=\footnotesize] at (axis cs:3.4,60) {$L$-край};
\end{axis}
\end{tikzpicture}
\caption{}
\end{subfigure}\hfill
\begin{subfigure}[b]{0.48\textwidth}\centering
\begin{tikzpicture}[scale=0.9]
  \foreach \lab/\y in {K/0, L/2.2, M/3.3, N/3.9}{\draw[line width=0.8pt] (0,\y) -- (4.6,\y); \node[slbl, right] at (4.6,\y) {$\lab$};}
  \draw[arr, fmRed] (0.5,2.2) -- (0.5,0.05) node[pos=0.5, left, slbl]{$K_\alpha$};
  \draw[arr, fmRed] (1.2,3.3) -- (1.2,0.05) node[pos=0.6, left, slbl]{$K_\beta$};
  \draw[arr, fmBlue] (2.4,3.3) -- (2.4,2.25) node[pos=0.5, left, slbl]{$L_\alpha$};
  \draw[arr, fmBlue] (3.1,3.9) -- (3.1,2.25) node[pos=0.4, right, slbl]{$L_\beta$};
  \draw[arr, fmGreen] (4.0,3.9) -- (4.0,3.35) node[pos=0.5, right, slbl]{$M_\alpha$};
\end{tikzpicture}
\caption{}
\end{subfigure}
\caption{Рентгеновские спектры многоэлектронного атома: \textit{а}~—
края поглощения (коэффициент поглощения скачком растёт, когда квант
способен выбить электрон из $K$- или $L$-слоя); \textit{б}~— схема
образования серий характеристического излучения}
\label{fig:xray}
\end{figure}

\paragraph{Порядок заполнения: правило Клечковского}
Поскольку энергия зависит от $n$ и $l$, при росте атомного номера уровни
заполняются не строго «по слоям». Обобщая опытные данные,
В.\,М.~Клечковский сформулировал правило (рис.~\ref{fig:klechkovsky}):
\begin{glavnoe}
Уровни заполняются в порядке возрастания суммы $n+l$, а при равных $n+l$~—
в порядке возрастания $n$: $1s,\,2s,\,2p,\,3s,\,3p,\,4s,\,3d,\,4p,\,5s,\,4d,\,5p,\,6s,\,4f,\,5d,\ldots$
\end{glavnoe}
Правило можно получить и теоретически из статистической модели атома
Томаса--Ферми. У него есть небольшие исключения (Cr, Cu, Pd и др.).

\begin{figure}[ht]
\centering
\begin{tikzpicture}[scale=0.82]
  \foreach \n in {1,...,7}{
    \foreach \l/\lt in {0/s,1/p,2/d,3/f}{
      \pgfmathtruncatemacro{\ok}{(\l<\n) ? 1 : 0}
      \ifnum\ok=1 \node[draw=fmBlue, fill=fmBlue!8, rounded corners=1pt, minimum width=8mm,
        minimum height=5mm, font=\footnotesize] at (\l*1.25,-\n*0.8) {$\n\lt$};\fi}}
  % диагонали n + l = s
  \foreach \s in {1,...,8}{
    \pgfmathtruncatemacro{\lmax}{min(3,ceil(\s/2)-1)}
    \pgfmathtruncatemacro{\lmin}{max(0,\s-7)}
    \draw[fmRed!75, -{Stealth[length=1.8mm]}, line width=0.7pt]
      ({\lmax*1.25+0.55},{-(\s-\lmax)*0.8+0.35}) -- ({\lmin*1.25-0.55},{-(\s-\lmin)*0.8-0.35});}
\end{tikzpicture}
\caption{Порядок заполнения электронных оболочек по правилу Клечковского:
вдоль каждой красной стрелки сумма $n+l$ постоянна; оболочки заполняются
стрелка за стрелкой сверху вниз, а вдоль стрелки~— справа налево}
\label{fig:klechkovsky}
\end{figure}

\paragraph{Правило Хунда}
Внутри одной оболочки ($p$, $d$, $f$) электроны располагаются по ячейкам
так, чтобы число \emph{параллельных} спинов было наибольшим: сначала по
одному в каждую ячейку, и лишь затем парами (рис.~\ref{fig:period2}).
Нарушение этого правила обошлось бы дорого: энергия атома углерода была
бы больше на $1{,}26\eV$, азота~— на $2{,}39\eV$, кислорода~— на
$1{,}96\eV$.

\begin{figure}[ht]
\centering
\begin{tikzpicture}[scale=0.62]
  \newcommand{\cellbox}[3]{% x, y, содержимое: 0 пусто, 1 вверх, 2 пара
    \draw[fill=white] (#1,#2) rectangle ++(0.8,0.8);
    \ifcase#3 \or \draw[-{Stealth[length=1.4mm]}, line width=0.7pt, fmRed] (#1+0.4,#2+0.12) -- (#1+0.4,#2+0.68);
    \or \draw[-{Stealth[length=1.4mm]}, line width=0.7pt, fmRed] (#1+0.27,#2+0.12) -- (#1+0.27,#2+0.68);
        \draw[-{Stealth[length=1.4mm]}, line width=0.7pt, fmBlue] (#1+0.53,#2+0.68) -- (#1+0.53,#2+0.12);\fi}
  \foreach \el/\a/\b/\c/\d/\e [count=\i from 0] in {Li/2/1/0/0/0, Be/2/2/0/0/0, B/2/2/1/0/0, C/2/2/1/1/0,
       N/2/2/1/1/1, O/2/2/2/1/1, F/2/2/2/2/1, Ne/2/2/2/2/2}{
    \pgfmathsetmacro{\xx}{\i*0}
    \node[lbl, anchor=east] at (-0.3,-\i*1.1+0.4) {\el};
    \cellbox{0}{-\i*1.1}{\a}
    \cellbox{1.4}{-\i*1.1}{\b}
    \cellbox{2.8}{-\i*1.1}{\c}\cellbox{3.6}{-\i*1.1}{\d}\cellbox{4.4}{-\i*1.1}{\e}
  }
  \node[slbl] at (0.4,1.3) {$1s$}; \node[slbl] at (1.8,1.3) {$2s$}; \node[slbl] at (3.6,1.3) {$2p$};
\end{tikzpicture}
\caption{Заполнение квантовых ячеек у элементов второго периода (принцип
Паули и правило Хунда)}
\label{fig:period2}
\end{figure}

\paragraph{Валентность}
Валентность обычно равна числу неспаренных электронов. Она может
меняться при переходе атома в возбуждённое состояние: электрон из
заполненной ячейки переходит в соседнюю свободную, и неспаренных
электронов становится больше. Так, бериллий в основном состоянии
$1s^22s^2$ нульвалентен, а в возбуждённом $1s^22s2p$~— двухвалентен;
бор из одновалентного становится трёхвалентным, углерод из
двухвалентного~— четырёхвалентным. Энергию на возбуждение даёт сама
химическая реакция.

\subsubsection{Периодическая система}

Второй период начинается с лития (электрон $2s$ сверх оболочки гелия)
и кончается неоном, у которого слой $n=2$ заполнен полностью: облако
сферически симметрично, оболочка очень устойчива. Элементы одной
вертикальной подгруппы имеют одинаковое строение внешних оболочек~—
отсюда \term{элементы-аналоги}. Третий период (Na\ldots Ar) тоже содержит
восемь элементов: по правилу Клечковского уровень $3d$ лежит выше $4s$, и
$3d$-электроны появляются только в четвёртом периоде~— после калия и
кальция.

\begin{opredelenie}
\term{Переходные металлы}~— элементы, у которых заполняется внутренняя
$(n-1)d$-оболочка при заполненной внешней $ns$-оболочке.
\end{opredelenie}

Особенности заполнения:
\begin{itemize}
\item У \emph{хрома} один из электронов $4s$ переходит в $3d$
(конфигурация $3d^54s^1$); у марганца аномалия исчезает. У \emph{меди}
аналогично $3d^{10}4s^1$: $3d$-оболочка заполнена. Хотя у меди, как у
щелочных металлов, один внешний $s$-электрон, на её свойства влияет
заполненная $d$-оболочка, и медь относят к другой подгруппе I группы.
\item У \emph{палладия} оба $5s$-электрона переходят на $4d$: изолированный
атом Pd~— $4d^{10}$ без $5s$-электронов.
\item В шестом периоде у \emph{лантана} появляется $5d$-электрон (вопреки
правилу Клечковского), но уже у церия начинает заполняться $4f$; 14
элементов с заполняющейся $4f$-оболочкой~— \term{лантаноиды}. В седьмом
периоде заполнение $5f$ даёт \term{актиноиды}.
\end{itemize}

\begin{glavnoe}
Недостроенная $d$-оболочка определяет многие свойства переходных металлов.
Из-за близости уровней $ns$ и $(n-1)d$ в связи участвуют электроны обоих
типов: отсюда переменная валентность, высокая прочность межатомной связи
(среди переходных металлов и их сплавов~— самые тугоплавкие и твёрдые
вещества с наибольшими модулями упругости) и склонность к полиморфизму.
\end{glavnoe}

\paragraph{Классификация элементов}
По способности образовывать прочные связи элементы делятся на химически
инертные (благородные газы He, Ne, Ar, Kr, Xe, Rn) и активные; активные~—
на металлы и неметаллы (табл.~\ref{tab:elements}). Металлов среди 104
элементов таблицы Менделеева (на момент написания оригинала)~— 85, т.\,е.
более 80\,\%.

\begin{table}[ht]
\centering
\caption{Классификация элементов}
\label{tab:elements}
\small
\begin{tabularx}{\textwidth}{@{}>{\raggedright\arraybackslash}p{38mm}>{\raggedright\arraybackslash}X@{}}
\toprule
Группа & Состав \\
\midrule
Истинные металлы (69) & щелочные (Li\ldots Fr), щёлочноземельные (Be\ldots Ra),
  переходные (Sc\ldots Ni, Y\ldots Pd, Hf\ldots Pt, La, Ac, Ku), лантаноиды (Ce\ldots Lu),
  актиноиды (Th\ldots Lr), металлы группы меди (Cu, Ag, Au) \\
Металлы подгруппы B (16) & B, Al, Si, Ga, Ge, As, Tl, Pb, Bi, Po, Zn, Cd, In, Sn, Sb, Hg~—
  свойства постепенно меняются от металлических к металлоидным \\
Неметаллы (13) & H, C, N, O, F, P, S, Cl, Se, Br, Te, I, At \\
Инертные (6) & He, Ne, Ar, Kr, Xe, Rn \\
\bottomrule
\end{tabularx}
\end{table}

\paragraph{Периодичность свойств}
Периодичность строения внешних оболочек порождает периодичность свойств.
Ярче всего это видно на первых потенциалах ионизации
(рис.~\ref{fig:ion-Z}): максимумы~— у благородных газов (самые
устойчивые оболочки), минимумы~— у щелочных металлов, начинающих новый
период.

\begin{figure}[ht]
\centering
\begin{tikzpicture}
\begin{axis}[width=0.95\textwidth, height=5.6cm, xmin=0, xmax=57, ymin=0, ymax=26,
  xlabel={атомный номер $Z$}, ylabel={$J_1$, эВ}, grid=major, grid style={fmGray!15}]
  \addplot[fmBlue, thick, mark=*, mark size=1.1pt] coordinates {
  (1,13.60)(2,24.59)(3,5.39)(4,9.32)(5,8.30)(6,11.26)(7,14.53)(8,13.62)(9,17.42)(10,21.56)
  (11,5.14)(12,7.65)(13,5.99)(14,8.15)(15,10.49)(16,10.36)(17,12.97)(18,15.76)
  (19,4.34)(20,6.11)(21,6.56)(22,6.83)(23,6.75)(24,6.77)(25,7.43)(26,7.90)(27,7.88)(28,7.64)(29,7.73)(30,9.39)
  (31,6.00)(32,7.90)(33,9.79)(34,9.75)(35,11.81)(36,14.00)
  (37,4.18)(38,5.69)(39,6.22)(40,6.63)(41,6.76)(42,7.09)(43,7.28)(44,7.36)(45,7.46)(46,8.34)(47,7.58)(48,8.99)
  (49,5.79)(50,7.34)(51,8.61)(52,9.01)(53,10.45)(54,12.13)(55,3.89)(56,5.21)};
  \node[font=\scriptsize, above] at (axis cs:2,24.59) {He};
  \node[font=\scriptsize, above] at (axis cs:10,21.56) {Ne};
  \node[font=\scriptsize, above] at (axis cs:18,15.76) {Ar};
  \node[font=\scriptsize, above] at (axis cs:36,14.00) {Kr};
  \node[font=\scriptsize, above] at (axis cs:54,12.13) {Xe};
  \node[font=\scriptsize, below] at (axis cs:3,5.39) {Li};
  \node[font=\scriptsize, below] at (axis cs:11,5.14) {Na};
  \node[font=\scriptsize, below] at (axis cs:19,4.34) {K};
  \node[font=\scriptsize, below] at (axis cs:37,4.18) {Rb};
  \node[font=\scriptsize, below] at (axis cs:55,3.89) {Cs};
\end{axis}
\end{tikzpicture}
\caption{Первые потенциалы ионизации элементов (справочные данные)}
\label{fig:ion-Z}
\end{figure}

\paragraph{Размеры атомов}
Размер изолированного атома точно определить нельзя: плотность облака
убывает с расстоянием медленно и обращается в нуль лишь на бесконечности.
В кристалле же за диаметр атома естественно принять кратчайшее
межатомное расстояние, которое рентгенография даёт очень точно. Для
металлов удобно пользоваться \term{атомным объёмом} $V_a=A/\rho$ ($A$~—
атомная масса, $\rho$~— плотность)~— объёмом одного моля; поделив его на
число Авогадро, получим объём на атом (рис.~\ref{fig:atvol}).

\begin{figure}[ht]
\centering
\begin{tikzpicture}
\begin{axis}[width=0.95\textwidth, height=5.6cm, xmin=0, xmax=84, ymin=0, ymax=75,
  xlabel={атомный номер $Z$}, ylabel={$V_a$, см$^3$/моль}, grid=major, grid style={fmGray!15}]
  \addplot[fmBlue, thick, mark=*, mark size=1.1pt] coordinates {
  (3,13.0)(4,4.9)(5,4.6)(6,3.4)
  (11,23.7)(12,14.0)(13,10.0)(14,12.1)
  (19,45.5)(20,26.2)(21,15.0)(22,10.6)(23,8.3)(24,7.2)(25,7.4)(26,7.1)(27,6.7)(28,6.6)(29,7.1)(30,9.2)(31,11.8)(32,13.6)
  (37,55.8)(38,33.9)(39,19.9)(40,14.0)(41,10.8)(42,9.4)(44,8.2)(45,8.3)(46,8.9)(47,10.3)(48,13.0)(49,15.7)(50,16.3)
  (55,70.0)(56,38.2)(57,22.4)(72,13.4)(73,10.9)(74,9.5)(75,8.9)(76,8.4)(77,8.5)(78,9.1)(79,10.2)(80,14.8)(81,17.2)(82,18.3)};
  \node[font=\scriptsize, above] at (axis cs:3,13.0) {Li};
  \node[font=\scriptsize, above] at (axis cs:11,23.7) {Na};
  \node[font=\scriptsize, above] at (axis cs:19,45.5) {K};
  \node[font=\scriptsize, above] at (axis cs:37,55.8) {Rb};
  \node[font=\scriptsize, above] at (axis cs:55,70.0) {Cs};
  \node[font=\scriptsize, below] at (axis cs:47,10.3) {Ag};
  \node[font=\scriptsize, below] at (axis cs:79,10.2) {Au};
  \node[font=\scriptsize, below] at (axis cs:28,6.6) {Ni};
  \draw[fmGray, dashed] (axis cs:58,2) -- (axis cs:71,2) node[midway, above, font=\scriptsize]{лантаноиды (не показаны)};
\end{axis}
\end{tikzpicture}
\caption{Атомные объёмы элементов в твёрдом состоянии (справочные данные):
максимумы~— у щелочных металлов, минимумы~— у переходных металлов.
Из-за лантаноидного сжатия золото по размеру не больше серебра}
\label{fig:atvol}
\end{figure}

\begin{itemize}
\item Наибольшие атомы~— у щелочных металлов: новый $s$-электрон
появляется сверх оболочки благородного газа.
\item Заполнение \emph{внутренних} $d$-оболочек при росте заряда ядра
стягивает внешнюю $s$-оболочку, и атомы уменьшаются. Наименьшие атомы и
межатомные расстояния~— у переходных металлов; у них же самые высокие
температуры плавления.
\item В шестом периоде перед $5d$-металлами заполняется ещё более глубокая
$4f$-оболочка (14 электронов лантаноидов). Это так сильно стягивает атомы,
что золото оказывается не больше серебра, хотя у него на 32 электрона
больше,~— \term{лантаноидное сжатие}. После него идут самые тугоплавкие
металлы W и Ta ($\Tpl\approx3400$ и $3000$\,°C).
\end{itemize}

Если построить зависимость обратной температуры плавления $1/\Tpl$ от
атомного номера, она окажется подобной зависимости атомного объёма от
$Z$; так же ведут себя коэффициент теплового расширения и сжимаемость.
Все эти величины отражают одно и то же~— прочность межатомной связи.

\begin{itogi}
\begin{itemize}
\item Частицы обладают волновыми свойствами, $\lambda=h/m\upsilon$;
состояние описывается волновой функцией, $|\psi|^2$~— плотность
вероятности.
\item Атом водорода: $E_n=-13{,}6\eV/n^2$; квантовые числа $n$, $l$,
$m_l$, $m_s$; $s$-облака сферические, $p$~— гантели.
\item Многоэлектронные атомы: слои $K,L,M,\ldots$, принцип Паули
($2n^2$ на слой), энергия зависит от $n$ и $l$; заполнение по правилу
$n+l$ (Клечковского), внутри оболочки~— по правилу Хунда.
\item Периодичность внешних оболочек $\Rightarrow$ периодичность
потенциалов ионизации, размеров атомов, $\Tpl$. Переходные металлы~—
малые атомы и прочная связь; лантаноидное сжатие.
\end{itemize}
\end{itogi}

\begin{proverka}
\item Чему равна длина волны де Бройля нейтрона с энергией $0{,}025\eV$
(тепловой нейтрон)?
\item Сколько электронов вмещает слой $M$? Оболочка $4f$?
\item Почему у лития второй потенциал ионизации в 14 раз больше первого?
\item В каком порядке заполняются уровни $3d$, $4s$, $4p$ и почему?
\item Почему атом золота не больше атома серебра?
\end{proverka}
